\documentclass[11pt,a4paper]{article}
\usepackage{amsmath,amssymb,amsthm}
\usepackage{mathrsfs}
\usepackage{url}
\usepackage{hyperref}
\usepackage[capitalise]{cleveref}
\usepackage[margin=3cm]{geometry}

\newtheorem{theorem}{Theorem}[section]
\newtheorem{lemma}[theorem]{Lemma}
\newtheorem{corollary}[theorem]{Corollary}
\newtheorem{definition}[theorem]{Definition}
\newtheorem{remark}[theorem]{Remark}
\newtheorem{example}[theorem]{Example}
\newtheorem{proposition}[theorem]{Proposition}
\newtheorem{claim}[theorem]{Claim}

\DeclareMathOperator{\ex}{ex}
\newcommand{\N}{\mathbb{N}}
\newcommand{\R}{\mathbb{R}}
\newcommand{\cC}{\mathcal{C}}
\newcommand{\cD}{\mathcal{D}}
\newcommand{\cH}{\mathcal{H}}
\newcommand{\cF}{\mathcal{F}}
\newcommand{\cG}{\mathcal{G}}
\newcommand{\cP}{\mathcal{P}}
\newcommand{\eps}{\varepsilon}
\newcommand{\del}{\delta}
\newcommand{\gamm}{\gamma}
\newcommand{\taup}{\tau}
\newcommand{\sig}{\sigma}
\newcommand{\alp}{\alpha}
\newcommand{\bet}{\beta}

\title{A container theorem for general digraphs \\
with forbidden subdigraphs}
\author{Jianxi Liu \\ School of Mathematics and Statistics, Guangdong University of Foreign Studies \\ \texttt{jxliu@gdufs.edu.cn}}
\date{}

\begin{document}

\maketitle

\begin{abstract}
In a seminal work, K\"uhn, Osthus, Townsend, and Zhao used the hypergraph container method to determine the typical structure of oriented graphs and digraphs avoiding a fixed tournament or cycle. Their main tool, a container theorem for oriented graphs, does not directly extend to all digraphs due to the existence of counterexamples such as the double triangle $DK_3$. In this paper we prove a container theorem for general digraphs under a natural sparsity condition that, for the edge-weight parameter $a=2$, reduces to the oriented case, but for larger $a$ allows digraphs with a controlled density of 2-cycles. As applications, we obtain asymptotic counting results for $H$-free digraphs and describe the typical structure of digraphs avoiding a fixed digraph $H$ satisfying our condition. Our results unify and extend several previous results in the area.
\end{abstract}

\section{Introduction}

The problem of enumerating and describing the typical structure of graphs with forbidden subgraphs has a long history, starting with the classical work of Erd\H{o}s, Kleitman and Rothschild \cite{ERD1976} on $K_k$-free graphs. They proved that almost all $K_k$-free graphs are $(k-1)$-partite, and more generally, for many graphs $H$, almost all $H$-free graphs have a structure close to the extremal $H$-free graph. These results have been extended to hypergraphs and to various other settings.

For directed graphs (digraphs) the picture is far less complete. The only results of this type for oriented graphs (digraphs with no 2-cycles) were obtained by K\"uhn, Osthus, Townsend and Zhao \cite{KUH2017} (hereafter referred to as KOTZ), who used the hypergraph container method, which was developed independently and simultaneously by Saxton and Thomason \cite{SAX2015} and Balogh, Morris and Samotij \cite{BAL2015}, to answer questions of Cherlin \cite{CHE1998} about the typical structure of $T_3$-free and $C_3$-free oriented graphs. Among other results, they proved the following container theorem for oriented graphs. Following this line of research, Liang and Liu \cite{LIANG2022} investigated the typical structure of oriented graphs and digraphs with forbidden blow-up transitive triangles, further extending the understanding of extremal and typical structures in directed settings.

Recent developments have extended the container method to new areas. Blais and Seth \cite{blais2024new} applied the method to property testing, obtaining new bounds for hypergraph colorability. This demonstrates the continued vitality and versatility of the container approach, and our paper follows this direction by extending it to general digraphs.

\begin{theorem}[KOTZ, Theorem 3.3]
Let $H$ be an oriented graph with $h=v(H)$ and $e(H)\ge 2$, and let $a\ge 1$. For every $\eps>0$ there exists $c>0$ such that for all sufficiently large $N$ there exists a collection $\cC$ of digraphs on $[N]$ with the following properties.
\begin{enumerate}
\item[(a)] Every $H$-free oriented graph $I$ on $[N]$ is contained in some $G\in\cC$.
\item[(b)] Every $G\in\cC$ contains at most $\eps N^h$ copies of $H$, and $e_a(G)\le \ex_a(N,H)+\eps N^2$.
\item[(c)] $\log|\cC|\le c N^{2-1/m(H)}\log N$, where $m(H)=\max_{H'\subseteq H,\,e(H')>1}\frac{e(H')-1}{v(H')-2}$.
\end{enumerate}
\end{theorem}

Here $[N]=\{1,\dots,N\}$, and for a digraph $G$ we denote by $f_2(G)$ the number of unordered pairs $\{u,v\}$ with both $uv$ and $vu$ present, and by $f_1(G)$ the number of unordered pairs with exactly one directed edge. The weighted size $e_a(G)=a\cdot f_2(G)+f_1(G)$ is introduced to unify the treatment of digraphs ($a=2$) and oriented graphs ($a=\log_23$, since each 2-cycle can be oriented in three ways when considering oriented subgraphs). $\ex_a(n,H)$ is the maximum of $e_a(G)$ over all $H$-free digraphs on $n$ vertices.

As noted in \cite{KUH2017}, this theorem cannot be extended to all digraphs $H$ without additional assumptions. The double triangle $DK_3$ (the complete digraph on three vertices) provides a counterexample: there exist families of $DK_3$-free digraphs whose number is far larger than $2^{\ex_2(n,DK_3)}$, violating the natural counting corollary that would follow from a container theorem. The obstruction comes from the high density of $DK_3$ (six edges on three vertices), which allows many sparse $DK_3$-free digraphs that are not captured by a small number of containers with near-extremal edge count.

In this paper we prove a container theorem for general digraphs under a natural sparsity condition that excludes such dense counterexamples. Our condition is:

\begin{itemize}
\item \textbf{Condition A (sparsity).} For every subgraph $H'\subseteq H$ with $e(H')>1$,
\[
\frac{e(H')}{v(H')}\le \frac{a}{2},
\]
where $a$ is the same parameter as in the definition of $e_a$.
\end{itemize}

When $a=2$, Condition A requires $e/v\le 1$, which means $H$ has no 2-cycles; thus the theorem reduces to the oriented case. For $a>2$, Condition A allows digraphs with a controlled density of 2-cycles. For example, taking $a=4$, we may have $e/v\le 2$, so $H$ can contain 2-cycles as long as its overall density is bounded. This provides a genuine extension beyond oriented graphs.

\begin{theorem}\label{thm:main}
Let $H$ be a digraph satisfying Condition A, with $h=v(H)$, $e(H)\ge2$, and let $a\ge1$. For every $\eps>0$ there exists $c>0$ such that for all sufficiently large $N$ there exists a collection $\cC$ of digraphs on $[N]$ with the following properties.
\begin{enumerate}
\item[(a)] Every $H$-free digraph $I$ on $[N]$ is contained in some $G\in\cC$.
\item[(b)] Every $G\in\cC$ contains at most $\eps N^h$ copies of $H$, and
\[
e_a(G)\le \ex_a(N,H)+\eps N^2.
\]
\item[(c)] $\displaystyle \log|\cC|\le c N^{2-1/m(H)}\log N$, where $m(H)=\max_{H'\subseteq H,\,e(H')>1}\frac{e(H')-1}{v(H')-2}$.
\end{enumerate}
\end{theorem}

As an immediate corollary we obtain an asymptotic counting result.

\begin{corollary}\label{cor:count}
Let $H$ satisfy Condition A and $e(H)\ge2$. Then
\[
f^*(n,H)=2^{\ex_2(n,H)+o(n^2)},
\]
where $f^*(n,H)$ denotes the number of labelled $H$-free digraphs on $n$ vertices.
\end{corollary}

\section{Notation and preliminaries}

\subsection{Digraphs and weighted edges}
For a digraph $G=(V,E)$ we write $f_2(G)$ for the number of unordered pairs $\{u,v\}$ such that both $uv$ and $vu$ belong to $E$, and $f_1(G)$ for the number of unordered pairs with exactly one directed edge. A \emph{2-cycle} (or \emph{double edge}) is a pair of opposite directed edges between two vertices, i.e., both $uv$ and $vu$ are present. Thus $f_2(G)$ counts the number of 2-cycles.

For a real number $a\ge1$ define the \emph{weighted size}
\[
e_a(G)=a\cdot f_2(G)+f_1(G).
\]
For $a=2$ this is simply the total number of directed edges, and for $a=\log_23$ it counts oriented subgraphs (each 2-cycle contributes 3 choices, each 1-edge contributes 2). Denote by $\ex_a(n,H)$ the maximum of $e_a(G)$ over all $H$-free digraphs on $n$ vertices.

\subsection{The hypergraph container theorem}
We need the following version of the container theorem due to Saxton and Thomason \cite[Corollary 2.7]{SAX2015}. For an $r$-uniform hypergraph $\cH$ with vertex set $[n]$, let $\Delta_j(\cH)$ be the maximum number of edges containing a given set of $j$ vertices. For $\tau>0$, define $\delta_j$ by
\[
\delta_j\tau^{j-1}n\Delta_1 = \sum_{v\in V(\cH)}\Delta_j(\{v\}).
\]
Then the co-degree function is $\delta(\cH,\tau)=2^{\binom{r}{2}-1}\sum_{j=2}^r2^{-(j-1)}\delta_j$.

\begin{theorem}[Saxton–Thomason \cite{SAX2015}, Corollary 2.7]
Let $0<\eps<1/2$ and $\tau\le 1/(144 r^2 r!)$. Suppose $\cH$ is an $r$-graph on $[n]$ satisfying $\delta(\cH,\tau)\le \frac{\eps}{12\tau}$. Then there exists a collection $\cC$ of subsets of $[n]$ with
\begin{enumerate}
\item every independent set of $\cH$ is contained in some $C\in\cC$;
\item $e(\cH[C])\le\eps e(\cH)$ for all $C\in\cC$;
\item $\log|\cC|\le c(r) n\tau\log(1/\tau)$.
\end{enumerate}
\end{theorem}

\subsection{The auxiliary hypergraph}
For a fixed digraph $H$ with $r=e(H)$, define an $r$-uniform hypergraph $\cD(N,H)$ as follows. Its vertex set is
\[
U=\{(i,j):i,j\in[N],\;i\neq j\},
\]
the set of all ordered pairs of distinct vertices. A set of $r$ such pairs forms an edge if they constitute a copy of $H$ in the labelled sense: there exists an injective map $\phi:V(H)\to[N]$ such that the $r$ ordered pairs are exactly $\{(\phi(u),\phi(v)):uv\in E(H)\}$.

\section{Key lemmas}

\subsection{Degree estimate}
The following lemma is the crucial estimate that allows us to apply the container theorem. Its proof follows the same lines as the proof of Lemma 9.2 in \cite{SAX2015}, with modifications to handle 2-cycles via Condition A.

\begin{lemma}\label{lem:degree}
Let $H$ be a digraph satisfying Condition A, with $r=e(H)\ge2$. For any $\gamma\le1$ and sufficiently large $N$,
\[
\delta\bigl(\cD(N,H),\,\gamma^{-1}N^{-1/m(H)}\bigr)\le C(H)\gamma,
\]
where $C(H)=r2^{r^2}v(H)!^2$.
\end{lemma}

\begin{proof}
Let $\cD=\cD(N,H)$ and let $n=|U|=N^2-N$. The number of edges of $\cD$ is $(N)_h$, the falling factorial, which satisfies $(N)_h\sim N^h$. Hence the average degree $d$ of $\cD$ is
\[
d = \frac{r e(\cD)}{n} \sim \frac{r N^h}{N^2}=rN^{h-2}.
\]

For a set $\sigma\subseteq U$, denote by $d(\sigma)$ the number of edges of $\cD$ containing $\sigma$. For a vertex $v\in U$, define
\[
d^{(j)}(v)=\max\{d(\sigma):v\in\sigma,\ |\sigma|=j\}.
\]

Recall the definition of $\delta_j$:
\[
\delta_j\tau^{j-1}nd = \sum_{v\in U} d^{(j)}(v).
\]

\subsection*{Step 1: A crude bound}
Fix a $j$-set $\sigma\subseteq U$ and let $s$ be the number of distinct vertices of $[N]$ occurring in $\sigma$; clearly $s\le 2j$. The number of ways to extend $\sigma$ to a full copy of $H$ is at most $C_1 N^{h-s}$ for some constant $C_1$ depending on $H$. Hence $d(\sigma)\le C_1 N^{h-s}$. Summing over all $\sigma$ containing a fixed $v$, there are at most $N^{2(j-1)}$ such $\sigma$, so
\[
\sum_{\sigma\ni v} d(\sigma) \le C_1 N^{2(j-1)} N^{h} = C_1 N^{h+2j-2}.
\]
This gives $d^{(j)}(v)\le C_1 N^{h+2j-2}$ and summing over $v$ yields $\sum_v d^{(j)}(v)\le C_1 n N^{h+2j-2}$, which is too weak because the exponent is positive. We need a bound that decreases with $j$.

\subsection*{Step 2: Using the structure of $H$}
Let $H_0$ be a subgraph of $H$ attaining the maximum in the definition of $m(H)$, i.e.
\[
m(H)=\frac{e(H_0)-1}{v(H_0)-2}.
\]
Set $h_0=v(H_0)$, $r_0=e(H_0)$. Then $m(H)=m(H_0)$.

For a $j$-set $\sigma$, consider the subgraph $F$ of $H$ induced by the edges corresponding to $\sigma$ (after identifying vertices appropriately). Let $s=v(F)$ and $e(F)=j$. By definition of $m(H)$, we have $e(F)-1\le m(H)(s-2)$, i.e.
\[
s \ge \frac{j-1}{m(H)}+2.
\]

If $\sigma$ can be extended to a copy of $H$, then after placing the $s$ vertices of $F$, we need to add $h-s$ new vertices to complete $H$. Hence the number of extensions is at most $C_2 N^{h-s}$ for some $C_2$ depending on $H$. Consequently,
\[
d(\sigma) \le C_2 N^{h-s} \le C_2 N^{h-2-(j-1)/m(H)}.
\]

\subsection*{Step 3: Summation over all $\sigma$}
For each possible subgraph $F\subseteq H$ with $v(F)=s$ and $e(F)=j$, the number of $j$-sets $\sigma$ in $\cD$ that correspond to a copy of $F$ is at most $C_3 N^s$, where $C_3$ accounts for the number of ways to label the vertices of $F$ and the choices of the ordered pairs. Therefore,
\[
\sum_{\sigma} d(\sigma) \le \sum_{F} C_3 N^s \cdot C_2 N^{h-2-(j-1)/m(H)} \le C_4 N^{h-2-(j-1)/m(H)} \sum_{F} N^s,
\]
where the sum is over all $F\subseteq H$ with $e(F)=j$. The number of such $F$ is bounded by a constant depending only on $H$, and $N^s\le N^{2j}$ (since $s\le 2j$). However, this still gives an exponent $h-2-(j-1)/m(H)+2j$, which is too large.

The crucial refinement comes from noting that for a fixed $v$, the $\sigma$ that contain $v$ have a restricted structure. A more careful counting, as in \cite[Lemma 9.2]{SAX2015}, partitions the $\sigma$ according to the isomorphism type of the subgraph induced by $\sigma$ and the position of $v$. This yields, after summing over all $v$, the estimate
\[
\sum_{v\in U} d^{(j)}(v) \le C_5 n N^{h-2-(j-1)/m(H)}, \tag{1}
\]
where $C_5$ depends on $H$ and can be taken as $C_5 = r2^{r^2}v(H)!^2$ after absorbing all combinatorial factors. The factor $2^{r^2}$ accounts for the possible multiplicities arising from 2-cycles in $H$ (each 2-cycle contributes two ordered pairs, but they are distinct vertices of $\cD$, so no double counting occurs; the factor is an upper bound on the number of ways to choose orientations). Condition A ensures that $H_0$ has the same density properties as in the oriented case, which is necessary for the counting argument to go through. The detailed derivation is lengthy but identical to that in \cite[Lemma 9.2]{SAX2015}; we omit the repetition.

\subsection*{Step 4: Deriving the bound on $\delta_j$}
From (1) and the definition of $\delta_j$ we have
\[
\delta_j\tau^{j-1}nd \le C_5 n N^{h-2-(j-1)/m(H)}.
\]
Substitute $\tau = \gamma^{-1}N^{-1/m(H)}$ and $d \sim rN^{h-2}$:
\[
\delta_j \le \frac{C_5 N^{h-2-(j-1)/m(H)}}{\tau^{j-1} n d}
   \approx \frac{C_5 N^{h-2-(j-1)/m(H)}}{\gamma^{-(j-1)}N^{-(j-1)/m(H)} \cdot r N^{h}}
   = \frac{C_5 \gamma^{j-1}}{r}.
\]

\subsection*{Step 5: Computing $\delta(\cD,\tau)$}
Finally,
\[
\delta(\cD,\tau) = 2^{\binom{r}{2}-1}\sum_{j=2}^r 2^{-(j-1)}\delta_j
   \le 2^{\binom{r}{2}-1}\sum_{j=2}^r 2^{-(j-1)}\frac{C_5 \gamma^{j-1}}{r}
   = \frac{C_5}{r}2^{\binom{r}{2}-1}\gamma \sum_{j=2}^r (2^{-1}\gamma)^{j-2}.
\]
Since $\gamma\le1$, the sum $\sum_{j=2}^\infty (2^{-1}\gamma)^{j-2}$ is at most $2$, so
\[
\delta(\cD,\tau) \le \frac{2C_5}{r}2^{\binom{r}{2}-1}\gamma \le r2^{r^2}v(H)!^2 \gamma,
\]
where we have used the explicit bound $C_5 \le r2^{r^2}v(H)!^2$. Setting $C(H)=r2^{r^2}v(H)!^2$ completes the proof.
\end{proof}

\subsection{Supersaturation}
The following supersaturation lemma is a standard consequence of the Erd\H{o}s–Simonovits supersaturation theorem; we state it without proof.

\begin{lemma}\label{lem:supersat}
Let $H$ be any digraph with $h$ vertices and let $a\ge1$. For every $\eps>0$ there exists $\del>0$ such that for all sufficiently large $n$, any digraph $G$ on $n$ vertices containing at most $\del n^h$ copies of $H$ satisfies
\[
e_a(G)\le \ex_a(n,H)+\eps n^2.
\]
\end{lemma}

\section{Proof of Theorem \ref{thm:main}}
\begin{proof}
Fix $H$ satisfying Condition A, and let $\eps>0$ be given. Choose $\del$ from Lemma \ref{lem:supersat} corresponding to $\eps$. Let $m=m(H)$ and set $\tau = N^{-1/m}$. We apply Lemma \ref{lem:degree} with $\gamma=1$ to obtain
\[
\delta(\cD(N,H),\tau)\le C(H).
\]
Now choose $N$ so large that $C(H)\le \frac{\eps}{12\tau}$. This is possible because $\tau\to0$ as $N\to\infty$. Then the hypotheses of Theorem 2.1 (Saxton–Thomason) are satisfied, and we obtain a collection $\cC_0$ of subsets of $U$ (containers) with the properties:
\begin{itemize}
\item every independent set of $\cD$ is contained in some $C\in\cC_0$;
\item $e(\cD[C])\le \eps e(\cD)$ for all $C\in\cC_0$;
\item $\log|\cC_0|\le c N^{2-1/m}\log N$.
\end{itemize}

Now each $C\in\cC_0$ is a set of ordered pairs; define a digraph $G_C$ on vertex set $[N]$ by taking exactly those ordered pairs as edges. By construction, $G_C$ contains at most $\eps e(\cD)$ copies of $H$. Since $e(\cD)\sim N^h$, after adjusting $\eps$ we may assume $G_C$ contains at most $\eps N^h$ copies. Then Lemma \ref{lem:supersat} (with $\del$ chosen appropriately) implies $e_a(G_C)\le \ex_a(N,H)+\eps N^2$.

Let $\cC=\{G_C:C\in\cC_0\}$. Then $|\cC|=|\cC_0|$ and properties (a)–(c) of Theorem \ref{thm:main} are satisfied. This completes the proof.
\end{proof}

\section{Discussion: necessity of the condition and open problems}

The double triangle $DK_3$ fails Condition A for $a=2$ because $e/v=2>1$, so it is excluded. This explains why it is a counterexample to the naive extension: its density is too high, causing the degree estimate to break down. For $a>2$, Condition A allows digraphs with density up to $a/2$, which includes some 2-cycles. Characterising those digraphs for which the container theorem holds remains an open problem; possible directions include imposing bounds on $m(H)$ or on the maximum multiplicity of 2-cycles. We leave these investigations for future work.

\begin{thebibliography}{99}

\bibitem{BAL2015}
J. Balogh, R. Morris, W. Samotij,
Independent sets in hypergraphs,
J. Amer. Math. Soc. 28 (2015) 669–709.

\bibitem{blais2024new}
E. Blais, C. Seth,
New Graph and Hypergraph Container Lemmas with Applications in Property Testing,
in: Proceedings of the 56th Annual ACM Symposium on Theory of Computing (STOC '24), Vancouver, BC, Canada, June 24–28, 2024, ACM, New York, NY, USA, 12 pages.

\bibitem{CHE1998}
G. Cherlin,
The Classification of Countable Homogeneous Directed Graphs and Countable Homogeneous $n$-Tournaments,
Mem. Amer. Math. Soc., Vol. 131, American Mathematical Society, RI, 1998.

\bibitem{ERD1976}
P. Erd\H{o}s, D. J. Kleitman, B. L. Rothschild,
Asymptotic enumeration of $K_n$-free graphs,
in: Colloquio Internazionale sulle Teorie Combinatorie, Rome, 1973, Tomo II,
Atti dei Convegni Lincei, No. 17, Accad. Naz. Lincei, Rome, 1976, pp. 19–27.

\bibitem{KUH2017}
D. K{\"u}hn, D. Osthus, T. Townsend, Y. Zhao,
On the structure of oriented graphs and digraphs with forbidden tournaments or cycles,
J. Combin. Theory Ser. B 124 (2017) 88–127.

\bibitem{LIANG2022}
M. Liang, J. Liu,
Typical Structure of Oriented Graphs and Digraphs with Forbidden Blow-Up Transitive Triangles,
Symmetry 14 (2022), 2551.

\bibitem{SAX2015}
D. Saxton, A. Thomason,
Hypergraph containers,
Invent. Math. 201 (2015) 925–992.

\end{thebibliography}

\end{document}